% =====================================================================
%  Aula 2 — conteúdo das notas (incluído por notas.tex e pela apostila)
% =====================================================================
\aula{2}{O que é Aprendizado de Máquina?}[O que é ML?]

\begin{intuicao}[title={Objetivos da aula}]
Ao final desta aula você deve ser capaz de: (i) explicar o que significa um
algoritmo ``aprender'' e por que o desempenho é medido em casos novos;
(ii) distinguir problemas de previsão ($\hat y$) de problemas de inferência
($\hat\beta$); (iii) classificar uma tarefa como regressão, classificação ou
agrupamento; (iv) formular um problema econômico como tarefa de aprendizado de
máquina; e (v) reconhecer o fluxo básico de um projeto em \pkg{scikit-learn}.
\end{intuicao}

\section{Três perguntas}

Comecemos por três perguntas que um economista poderia receber numa mesma
semana. Quanto deve valer um apartamento de 70\,m² em Pinheiros? Um cliente que
pede um empréstimo vai pagá-lo? Que ``tipos'' de famílias existem nos dados de
uma pesquisa de orçamentos domésticos?

As três têm algo em comum: podem ser respondidas olhando para \emph{muitos casos
parecidos} e extraindo deles um padrão. Ninguém escreve à mão a regra que liga
área e localização ao preço; a regra é \emph{aprendida} a partir de transações
passadas. Mas as três também diferem no que queremos obter: na primeira, um
número; na segunda, uma categoria (paga ou não paga); na terceira, grupos que
ninguém rotulou antes. Essas diferenças dão nome aos principais tipos de
aprendizado de máquina, e reconhecê-las é o primeiro passo de qualquer projeto.

\section{O que significa ``aprender''?}

\begin{definicao}
Um programa \emph{aprende} com a experiência $E$, em relação a uma classe de
tarefas $T$ e a uma medida de desempenho $P$, se o seu desempenho nas tarefas
em $T$, medido por $P$, melhora com a experiência $E$ \parencite{mitchell1997}.
\end{definicao}

A definição é deliberadamente operacional: não fala de ``inteligência'', fala de
desempenho mensurável. No exemplo do crédito, a tarefa $T$ é prever se um novo
cliente ficará inadimplente; a experiência $E$ é o histórico de clientes passados,
com a informação de quem pagou; e o desempenho $P$ pode ser a proporção de
clientes \emph{novos} classificados corretamente.

A palavra \emph{novos} é o coração da definição. Um programa que simplesmente
guardasse a lista de clientes passados acertaria 100\% deles, e não teria
aprendido nada útil, porque o banco precisa decidir sobre quem ainda não está na
lista. Aprender é \emph{generalizar}; decorar não é aprender. A Aula 3 inteira
será dedicada a medir essa diferença com honestidade.

\subsection{A mesma ideia em linguagem econométrica}

Formalmente, supomos que o alvo $y$ e os atributos $\vx$ estão ligados por
\begin{equation}\label{eq:a02-modelo}
  y = f(\vx) + \varepsilon,
\end{equation}
em que $f$ é desconhecida e $\varepsilon$ reúne tudo o que $\vx$ não explica.
Dispomos de um conjunto de dados $\data = \{(\vx_i, y_i)\}_{i=1}^{n}$.
\emph{Aprender} é usar $\data$ para construir uma aproximação $\hf$;
\emph{usar o modelo} é calcular $\hat y = \hf(\vx_{\text{novo}})$ para casos
que não estavam em $\data$.

\begin{conexao}
Se impomos $f(\vx) = \vx\T\vbeta$ e escolhemos $\hvbeta$ minimizando a soma
dos quadrados dos resíduos, \cref{eq:a02-modelo} é o modelo de regressão linear
e o MQO é o algoritmo de aprendizado. Não há ruptura entre econometria e
aprendizado de máquina: há uma mudança de \emph{ênfase}, como veremos a seguir.
\end{conexao}

\subsection{Um pequeno dicionário}

As duas comunidades desenvolveram vocabulários próprios para as mesmas ideias.
A \cref{tab:a02-dicionario} reúne as traduções mais comuns; vale guardá-la, porque
a documentação das bibliotecas usa sempre o vocabulário da direita.

\begin{table}[ht]
  \centering\small
  \caption{Termos equivalentes em econometria e aprendizado de máquina.}
  \label{tab:a02-dicionario}
  \begin{tabular}{@{}lll@{}}
    \toprule
    Econometria & Aprendizado de máquina & Símbolo \\
    \midrule
    variável dependente      & alvo, rótulo (\emph{target}, \emph{label}) & $y$ \\
    variáveis explicativas   & atributos (\emph{features})                  & $\vx$ \\
    observação               & exemplo, instância                           & $(\vx_i,y_i)$ \\
    amostra                  & conjunto de treino                           & $\dtreino$ \\
    estimar                  & treinar, ajustar (\emph{fit})                & $\data\mapsto\hf$ \\
    parâmetros               & pesos (\emph{weights})                       & $\vbeta$, $\vw$ \\
    intercepto               & viés (\emph{bias})                           & $\beta_0$, $b$ \\
    valor ajustado           & previsão (\emph{prediction})                 & $\hat y$ \\
    \bottomrule
  \end{tabular}
\end{table}

\begin{atencao}
Em ML, ``viés'' (\emph{bias}) pode designar simplesmente o intercepto de um
modelo, sem relação com o viés de um estimador ($\E[\hat\theta]-\theta$). O
contexto decide. Na Aula 3 encontraremos o viés no segundo sentido, no
compromisso entre viés e variância.
\end{atencao}

\section{Estatística, econometria e aprendizado de máquina}

\subsection{Duas culturas}

\textcite{breiman2001cultures} descreveu duas culturas na análise de dados
(\cref{fig:a02-culturas}). Na \emph{cultura da modelagem de dados}, dominante
na estatística e na econometria, o pesquisador especifica um modelo para o
processo gerador dos dados (por exemplo, $y=\mX\vbeta+\ve$ com hipóteses sobre
$\ve$) e o interesse recai sobre os parâmetros: seus valores, seus erros-padrão,
sua interpretação. Na \emph{cultura algorítmica}, o processo gerador é tratado
como uma caixa-preta; escolhe-se o algoritmo que melhor prevê casos novos, e a
validação fora da amostra é o critério de qualidade.

\begin{figure}[ht]
  \centering
  \includegraphics[width=0.78\linewidth]{a02_duas_culturas}
  \caption{As duas culturas de \textcite{breiman2001cultures}. O que muda não é a
  matemática disponível, e sim o produto final que se deseja: $\hat\beta$ ou $\hat y$.}
  \label{fig:a02-culturas}
\end{figure}

\subsection{Problemas de \texorpdfstring{$\hat\beta$}{beta} e problemas de \texorpdfstring{$\hat y$}{y}}

\textcite{mullainathan2017} traduzem essa distinção para a economia aplicada.
Muitos trabalhos empíricos são \emph{problemas de $\hat\beta$}: qual o efeito de
um ano adicional de estudo sobre o salário? Um programa de transferência de renda
afeta a frequência escolar? Aqui o objetivo é entender um mecanismo, e a
dificuldade central é a \emph{identificação}: saber o que aconteceria se
mudássemos $x$, mantido o resto constante. Outros são \emph{problemas de
$\hat y$}: qual a probabilidade de um cliente ficar inadimplente? Qual será a
inflação do próximo mês? Aqui importa acertar casos ainda não observados, e a
dificuldade central é a \emph{generalização}.

O aprendizado de máquina foi construído para problemas de $\hat y$. Usá-lo para
responder perguntas de $\hat\beta$ sem cuidados adicionais é um erro comum, e
a \cref{fig:a02-prever} mostra por quê.

\begin{figure}[ht]
  \centering
  \includegraphics[width=0.85\linewidth]{a02_prever_inferir}
  \caption{Mil amostras de $n=100$ observações de $y = x_1 + x_2 + \varepsilon$,
  com $\operatorname{corr}(x_1,x_2)=0{,}98$. Em cada amostra estimamos o modelo por
  MQO. O coeficiente de $x_1$ varia muito entre amostras; a previsão no ponto
  $x_0 = (1,1)$ varia bem menos.}
  \label{fig:a02-prever}
\end{figure}

\begin{intuicao}
Quando $x_1$ e $x_2$ andam juntos, os dados informam bem o efeito
\emph{conjunto} $\beta_1+\beta_2$, que é o que a previsão em $x_0=(1,1)$ usa,
mas informam mal como esse efeito se divide entre $\beta_1$ e $\beta_2$.
Para prever, a divisão não importa; para interpretar, ela é tudo. Por isso um
modelo com ótimo desempenho preditivo pode ter coeficientes que mudam de sinal
de uma amostra para outra.
\end{intuicao}

\subsection{Guarda-chuva ou dança da chuva?}

\textcite{kleinberg2015} propõem uma imagem simples. Quem decide se leva o
guarda-chuva só precisa \emph{prever} se vai chover; o motivo da chuva é
irrelevante. Quem decide se faz a dança da chuva precisa saber se a dança
\emph{causa} chuva; nenhuma previsão responde a isso. Os autores mostram que
muitas decisões de política pública são do primeiro tipo, ou seja,
\emph{problemas de política preditiva}. Um exemplo do artigo: identificar
pacientes idosos que provavelmente morreriam no ano seguinte, para os quais
uma cirurgia de prótese de quadril traria pouco benefício. A decisão depende de
uma previsão, não de um efeito causal.

\begin{conexao}
Previsão e causalidade não são rivais. Métodos modernos de inferência causal
usam aprendizado de máquina em etapas puramente preditivas, por exemplo para
controlar de forma flexível muitas variáveis de confusão
\parencite{chernozhukov2018,athey2019}. Voltaremos a isso na Aula 6.
\end{conexao}

\section{Tipos de aprendizado}

A pergunta que organiza a taxonomia é: \emph{observamos o alvo $y$?}
(\cref{fig:a02-taxonomia}).

\begin{figure}[ht]
  \centering
  \includegraphics[width=0.82\linewidth]{a02_taxonomia}
  \caption{Tipos de aprendizado de máquina. Neste curso tratamos dos ramos
  supervisionados; o aprendizado por reforço fica fora do escopo.}
  \label{fig:a02-taxonomia}
\end{figure}

\begin{description}
  \item[Aprendizado supervisionado.] Observamos pares $(\vx_i, y_i)$: os dados
  ``supervisionam'' o algoritmo com respostas corretas. Se $y$ é contínuo
  (preço, salário, inflação), temos uma tarefa de \emph{regressão}; se $y$ é
  categórico (inadimplente ou não, setor de atividade), uma tarefa de
  \emph{classificação}.
  \item[Aprendizado não supervisionado.] Observamos apenas $\vx_i$. O objetivo é
  descobrir estrutura: grupos de observações parecidas (\emph{agrupamento}) ou
  poucas combinações de variáveis que resumem muitas (\emph{redução de
  dimensionalidade}, como num índice de condições financeiras construído a
  partir de dezenas de séries).
  \item[Aprendizado por reforço.] Um agente aprende a agir por tentativa e erro,
  recebendo recompensas do ambiente. É uma área ativa em economia computacional,
  mas exige ferramentas que estão além deste curso.
\end{description}

A \cref{fig:a02-tarefas} ilustra as três tarefas principais. Note a diferença
no terceiro painel: não há cores ``verdadeiras''. O algoritmo propõe três grupos
porque nós pedimos três; se pedíssemos cinco, ele proporia cinco. Interpretar e
validar os grupos é responsabilidade do analista.

\begin{figure}[ht]
  \centering
  \includegraphics[width=\linewidth]{a02_tarefas}
  \caption{Regressão (prever um número), classificação (prever uma categoria,
  com a fronteira de decisão tracejada) e agrupamento (propor grupos sem rótulos;
  os X marcam os centros encontrados pelo algoritmo $k$-médias). Dados simulados.}
  \label{fig:a02-tarefas}
\end{figure}

\subsection{Formular é escolher}

A tarefa não está escrita nos dados; é uma escolha do analista. Com a mesma
base de domicílios, podemos prever a renda em reais (regressão), prever se a
renda supera dois salários mínimos (classificação) ou procurar perfis de
domicílios sem alvo nenhum (agrupamento). Discretizar um alvo contínuo
simplifica a comunicação, mas descarta informação. O critério para escolher
é a \emph{decisão} que a previsão vai apoiar.

\section{Formulando problemas econômicos}

Antes de abrir o Python, vale preencher uma ficha com seis itens:

\begin{enumerate}[label=\arabic*.]
  \item \textbf{Decisão}: quem vai usar a previsão, e para quê?
  \item \textbf{Unidade de observação}: o que é uma linha da base?
  \item \textbf{Alvo $y$}: o que exatamente queremos prever, com definição precisa?
  \item \textbf{Atributos $\vx$}: que informação está disponível \emph{no momento
        em que a decisão é tomada}?
  \item \textbf{Tarefa}: regressão, classificação ou agrupamento?
  \item \textbf{Métrica}: como saberemos se o modelo é bom o suficiente para a decisão?
\end{enumerate}

\begin{exemplo}[title={Exemplo: crédito pessoal}]
\textbf{Decisão:} aprovar ou recusar um pedido de empréstimo.
\textbf{Unidade:} um pedido.
\textbf{Alvo:} atraso superior a 90 dias nos 12 meses seguintes (sim/não).
\textbf{Atributos:} renda declarada, idade, razão dívida/renda, histórico no
cadastro positivo, todos conhecidos na data do pedido.
\textbf{Tarefa:} classificação.
\textbf{Métrica:} os dois erros têm custos diferentes. Recusar um bom pagador
custa a margem do empréstimo; aprovar um mau pagador custa o principal. A
métrica deve refletir essa assimetria (Aula 3).
\end{exemplo}

\begin{atencao}
O item 4 é onde a maioria dos projetos tropeça. Se um atributo só é conhecido
\emph{depois} do evento que queremos prever (``número de cartas de cobrança
enviadas'' para prever inadimplência, ou a revisão final do PIB para prever o
próprio PIB), o modelo parecerá excelente nos dados históricos e será inútil na
prática. Esse fenômeno, o \emph{vazamento de dados}, é um dos temas da Aula 3.
\end{atencao}

\section{Fluxo de trabalho em Python}

Um projeto de aprendizado supervisionado segue, em linhas gerais, o fluxo da
\cref{fig:a02-fluxo}. As etapas em laranja (separar e avaliar) são o assunto da
Aula 3; a etapa em azul (treinar) será aberta por dentro nas Aulas 4 e 5.

\begin{figure}[ht]
  \centering
  \includegraphics[width=\linewidth]{a02_fluxo}
  \caption{Fluxo de um projeto de aprendizado supervisionado. O ciclo
  avaliar–ajustar pode se repetir várias vezes, mas sempre com dados de
  validação; o conjunto de teste é usado uma única vez, no final.}
  \label{fig:a02-fluxo}
\end{figure}

A biblioteca \pkg{scikit-learn} \parencite{pedregosa2011} padroniza esse fluxo:
todo modelo é um objeto com os métodos \py{fit} (aprender) e \py{predict}
(prever).

\begin{codigo}
from sklearn.model_selection import train_test_split
from sklearn.linear_model import LinearRegression
from sklearn.dummy import DummyRegressor
from sklearn.metrics import mean_squared_error

X_tr, X_te, y_tr, y_te = train_test_split(X, y, test_size=0.2, random_state=0)

for modelo in [DummyRegressor(strategy="mean"), LinearRegression()]:
    modelo.fit(X_tr, y_tr)
    rmse = mean_squared_error(y_te, modelo.predict(X_te)) ** 0.5
    print(type(modelo).__name__, round(rmse, 3))
\end{codigo}

Dois detalhes merecem destaque. Primeiro, o desempenho é calculado em
\py{X_te}, \py{y_te}, dados que o modelo não viu durante o \py{fit}.
Segundo, o laço compara o modelo com uma \emph{referência ingênua}, que prevê
sempre a média do treino.

\begin{atencao}
Um número de desempenho isolado não diz nada. Se 95\% dos clientes pagam em
dia, um ``modelo'' que declara que \emph{todos} pagam acerta 95\% dos casos e
não tem nenhuma utilidade. Compare sempre com uma referência ingênua: a média
(regressão), a classe mais frequente (classificação) ou o último valor
observado (séries temporais).
\end{atencao}

\section{Resumo e leituras}

\begin{itemize}[nosep]
  \item Aprender é melhorar o desempenho em \emph{casos novos} com a experiência.
  \item A regressão linear já é um algoritmo de aprendizado; o que muda no ML é a
        ênfase em $\hat y$ e na validação fora da amostra.
  \item Bons $\hat y$ não garantem $\hat\beta$ interpretáveis.
  \item Supervisionado (regressão, classificação) versus não supervisionado
        (agrupamento, redução de dimensão).
  \item Formular vem antes de modelar; os atributos precisam estar disponíveis no
        momento da decisão.
  \item \py{fit}, \py{predict} e uma referência ingênua, sempre.
\end{itemize}

\noindent\textbf{Leituras sugeridas.} \textcite{mullainathan2017} é a melhor
porta de entrada para economistas. \textcite[cap.~2]{james2021} apresenta a
formulação estatística com exemplos acessíveis. \textcite{breiman2001cultures}
e os comentários publicados junto ao artigo continuam atuais. Para um panorama
das aplicações em economia, \textcite{athey2019} e \textcite{varian2014}.
